% ════════════════════════════════════════════════════════════════════
% SECTION 9 — CONCLUSION
% ════════════════════════════════════════════════════════════════════
\section{Conclusion}\label{sec:conclusion}

We presented the first application of federated learning to solar
flare prediction for smart-grid protection, motivated by a structural
impasse in operational space weather: the high-fidelity observations
needed to train strong flare-forecasting models are held by national
agencies whose data-sovereignty commitments forbid pooling them. The
proposed framework partitions the SWAN-SF benchmark into six
simulated regional observatories under Dirichlet label skew, federates
MLP clients with FedAvg and FedProx under distribution-aware
aggregation, handles extreme class imbalance with per-client SMOTE and
a clamped server-aware focal loss, and converts converged models into
safety-critical detectors through recall-weighted $F_\beta$ threshold
optimisation.

The empirical outcome is two-sided and honest. On the side of promise:
FedProx retains 95.0\% of the centralised XGBoost ROC-AUC (0.930
against 0.978) and 77.6\% of its F1 (0.326 against 0.420) at a 2\%
false-alarm rate, converging within 50 communication rounds on
single-GPU compute---evidence that a jointly trained, globally
informed, sovereignty-preserving flare sentinel is technically within
reach, and that its interpretability layer recovers exactly the
magnetic parameters (flux-emergence proxies, current helicity,
Lorentz-force variability) that flare physics predicts. On the side of
caution: FedAvg collapses outright under the same conditions, the
training-time calibration of rebalanced federated data is broken
enough that threshold re-optimisation is mandatory rather than
optional, and the strongest claims (cryptographic privacy, LSTM and
SCAFFOLD evaluation, multi-seed statistics) remain future work. The
released framework---code, configuration history, and artefact
pipeline---is intended to make that follow-up work a matter of
execution rather than re-engineering.

For the smart-grid community the message is that resilience against
geomagnetic threats does not require choosing between data
sovereignty and forecast skill; for the machine-learning community it
is a demonstration that the cross-silo regime of the physical sciences
has failure modes---prior mismatch, layered imbalance handling,
optimiser--variance-reduction interactions---that the medical and
mobile FL literature does not fully anticipate, and that are worth
studying in their own right.
